\title{Algebra-Based Physics-1: Midterm 2}
\author{Dr. Jordan Hanson - Whittier College Dept. of Physics and Astronomy}
\date{\today}
\documentclass[10pt]{article}
\usepackage[margin=1.5cm]{geometry}
\usepackage{outlines}
\usepackage{graphicx}
\usepackage{amsmath}
\usepackage{hyperref}

\begin{document}
\twocolumn

\maketitle

\section{Unit 4: Rotational Kinematics}

\begin{enumerate}
\item Suppose a music record spins once every 1.33 seconds.  What are the rotations per minute (rpm) of the record?  What is the angular velocity, in radians per second?
\begin{itemize}
\item A: 30 rpm, 0.75 radians sec$^{-1}$
\item B: 45 rpm, 4.7 radians sec$^{-1}$
\item C: 60 rpm, 0.75 radians sec$^{-1}$
\item D: 120 rpm, 4.7 radians sec$^{-1}$
\end{itemize}
\item What is the centripetal acceleration of a coin at the edge of the record?
\begin{itemize}
\item A: 1.1 m s$^{-1}$
\item B: 1.1 m s$^{-2}$
\item C: 2.2 m s$^{-1}$
\item D: 2.2 m s$^{-2}$
\end{itemize}
\item Suppose a blood centrifuge is spinning at an angular velocity of $\omega$. In order to separate the contents in the vials, the centripetal acceleration needs to be increased by a factor of 100.  Which of the following actions will achieve this?
\begin{itemize}
\item A: $\omega \rightarrow 2\omega$
\item B: $\omega \rightarrow 3\omega$
\item C: $\omega \rightarrow 5\omega$
\item D: $\omega \rightarrow 10\omega$
\end{itemize}
\item Assume the angular velocity of the centrifuge is 200 rpm, and the radius is 12 cm.  If the solid contents of the vial in the centrifuge have a mass of 10 grams, what is the centripetal force?
\begin{itemize}
\item A: 0.13 N
\item B: 0.53 N
\item C: 0.83 N
\item D: 1.23 N
\end{itemize}
\item In lacrosse, a ball is thrown from a net on the end of a stick by rotating the stick and forearm about the elbow. If the angular velocity of the ball about the elbow joint is 30.0 rad/s and the ball is 1.40 m from the elbow joint, what is the velocity of the ball?
\begin{itemize}
\item A: 21 m s$^{-1}$
\item B: 30 m s$^{-1}$
\item C: 42 m s$^{-1}$
\item D: 50 m s$^{-1}$
\end{itemize}
\end{enumerate}

\section{Unit 5: Work and Energy}

\begin{enumerate}
\item Suppose a person pushes a box with a force $\vec{F} = 40\hat{i} - 30\hat{j}$ N of force, and the box moves $\Delta\vec{x} = 5\hat{i}$ m. What is the work done?  Wht is the angle between the force and displacement?
\begin{itemize}
\item A: 100 N m, 27 degrees
\item B: 100 N, 27 degrees
\item C: 200 N, 37 degrees
\item D: 200 N m, 37 degrees 
\end{itemize}
\item  Suppose a baseball player hits a baseball (0.15 kg) dropped in front of them with a baseball bat. (a) If the force during the hit is 1000 N, and the baseball is displaced 10 cm during the hit, what is the work done? (b) What is the change in kinetic energy, if we assume the baseball was initially stationary? (c) What is the final velocity of the ball? (d) If the ball leaves at a 45 degree angle, what is the range? \\ \vspace{2.5cm}
\item  Suppose we push a 40 kg box along a smooth floor with $\mu_k$ = 0.05 for a distance 40 m. (a) What is the work done? (b) If we perform this task in 30 seconds, what power is generated? \\ \vspace{2.5cm}
\item What is the cost of operating a 3.00 W electric clock for a year if the cost of electricity is 9 cents per kW h?
\begin{itemize}
\item A: \$6.70
\item B: \$13.30
\item C: \$1.67
\item D: \$2.33
\end{itemize}
\end{enumerate}

\section{Unit 6: Linear Momentum}

\begin{enumerate}
\item Two molecules collide and stick together, forming one larger molecule. Each molecule weighs $20 \times 10^{-25}$ kg.  One has a velocity of 350 m/s, and the other has a velocity of -350 m s$^{-1}$.  What is the final speed of the big new molecule?
\begin{itemize}
\item A: 350 m s$^{-1}$
\item B: -350 m s$^{-1}$
\item C: 0 m s$^{-1}$
\item D: 700 m s$^{-1}$
\end{itemize}
\item Assume the same two molecules approach each other with momentum vectors that are 45 degrees apart.  What is the final velocity after they stick together? \\ \vspace{2.5cm}
\item Suppose a child drives a bumper car head on into the side rail, which exerts a force of 4000 N on the car for 0.200 s. (a) What impulse is imparted by this force? (b) Find the final velocity of the bumper car if its initial velocity was 2.80 m/s and the car plus driver have a mass of 200 kg. You may neglect friction between the car and floor. \\ \vspace{2.5cm}
\item  Two identical billiard balls have a one-dimensional collision in which one is initially motionless. After the collision, the struck ball moves with the same velocity as the ball that was originally moving, while the ball that was originally moving stops.  This interaction is
\begin{itemize}
\item A: Totally inelastic
\item B: Inelastic
\item C: Elastic
\item D: Totally elastic
\end{itemize}
\item A 30,000-kg freight car is coasting at 0.850 m/s with negligible friction under a hopper that dumps 110,000 kg of scrap metal into it. (a) What is the final velocity of the loaded freight car? (b) How much kinetic energy is lost? \\ \vspace{2cm}
\end{enumerate}

\section{Unit 7: Statics}

\begin{enumerate}
\item A steel beam is left balancing on one support in the middle. A 40 kg pile of bricks rests 3 m from the origin, and a 60 kg sack of concrete lies 2 m from the origin.  The beam will:
\begin{itemize}
\item A: Remain motionless
\item B: Rotate towards the bricks
\item C: Rotate towards the concrete
\end{itemize}
\item A person carries a plank of wood 2.00 m long with one hand pushing down on it at one end with a force $\vec{F}_1$ and the other hand holding it up at 0.500 m from the end of the plank with force $\vec{F}_2$.  If the plank has a mass of 20.0 kg and its center of gravity is at the middle of the plank, what are the magnitudes of the forces $\vec{F}_1$ and $\vec{F}_2$? \textit{Hint: the net force and the net torque should both be zero (the two stability criteria).}\\ \vspace{3cm}
\end{enumerate}

\section{Unit 8: Rotational Dynamics}

\begin{enumerate}
\item Suppose we are pulling a rip chord to start a boat motor. We spin an internal gear shaped like a disc, and it has a moment of inertia $I = \frac{1}{2} M R^2$. (a) If the gear is 1 kg, and the radius is 8 cm, what is the moment of inertia in kg m$^2$? (b) What torque is required from us to spin the gear with $\alpha = 2$ rad s$^{-2}$? \\ \vspace{3cm}
\item To tighten the axle bolt on a bicycle, a mechanic presses one pedal forward, rotating it.  Let the location of the axle be the origin, and let the pedal position be $\vec{r} = -20\hat{i}+20\hat{j}$ cm.  (a) If the force on the pedal is $\vec{F} = 12\hat{i}$ N, what is the torque on the axle bolt? (b) What force vector is required to produce the opposite torque?
\end{enumerate}

\end{document}